\section{Data and evaluation protocol}
\label{sec:protocol}

\paragraph{Data.}
We use 1,223 full-disk AIA images from 2014 and 2015, about two per day. We deconvolve them with revised point-spread functions~\cite{hofmeister2025psf} and correct for exposure time and degradation. DEMs live on a $2048\times2048$ grid, with DEM pixel $(i,j)$ reading AIA pixel $(2i,2j)$. We cut each image's grid into 256 blocks and keep a random 64. The split is in time: 917 images for training (to 16 June 2015), 153 for validation and 153 for test (21 September to 31 December 2015, on 86 distinct days).

\paragraph{Shared test set.}
We use the evaluation protocol and scripts of DEMNet's authors and score both models on the same pixels against the same solver DEMs. Each test image has five targets, the solve of the clean image and four re-solves with Gaussian noise of $0.5\sigma$ added, and the input is the clean image for all five. The protocol scores each target on a random 64 of the image's 256 blocks, giving 48,960 blocks and 802M pixel-targets on Elastic Net. On BP 705M of these carry a solver DEM. The missing share at clean solves, 7.7\%, is close to the 8.5\% of pixels our own solver run had to relax, so the missing targets appear to be the infeasible pixels, and BP scores then describe feasible pixels only.

\paragraph{Brightness.}
A pixel is Bright when any channel reaches that channel's top-5\% cutoff and Quiet otherwise. Flare cores are pixels with some channel at 32 times its cutoff or more, which in 94\Ang{} means 100\,DN/s: 69,283 of the 705M, about one pixel in 10,000. They are the brightest pixels of flaring active regions rather than flare peaks alone. The strongest flare in the test window is close to X1, on 28 September 2015.

\paragraph{Measures.}
DEM MAE and MSE average over pixels and the 18 bins. EM error is the summed absolute error in total emission over the summed solver emission. W1 is the Wasserstein distance in dex between the two DEMs normalized as distributions over \logT. AIA MAE and MSE compare each DEM's predicted intensities with the observation in DN/s. These six are the headline measures. The per-pixel error is the squared error summed over the bins, and we report its percentiles because it is heavy-tailed: on an early model, one pixel in five million held 94\% of the mean training loss and reversed a model ranking. Flare-core emission and 94\Ang{} reproduced are ratios of sums over flare cores.

\paragraph{Baseline.}
The supervised baseline uses the DEMNet architecture~\cite{perezdiaz2026demnet}: 1.05M parameters, regression and classification heads over the DEM bins, and 140 epochs on solver DEMs from the same deconvolved images, with a resynthesis term of weight 0.01 from epoch 100. We use one checkpoint per solver from DEMNet's authors and evaluate both ourselves, on the same pixels and with the same code as our models. On BP the checkpoint scores DEM MSE 0.91 and AIA MSE 157.

\paragraph{Model selection.}
We chose widths, inputs and checkpoints on the training loss, which needs no solver DEM. During training we logged it on 131,072 pixels from the first four validation images, a sample too small for the rarest pixels that, with clamped pixels dropped, contains none of the inputs that fail. We therefore also scored the candidates on every pixel of the 153 validation images (\cref{sec:allpixels}), and that pass supports the final BP model. The Elastic Net width, epochs and treatment of clamped pixels rest on test scores, as does the rejection of the widened band in \cref{sec:infeasible}. The penalty weight came from a screen against solver DEMs on validation. \Cref{app:selection} gives the details.
